% [Lynn] Conclusion (rubric: 1 mark).
% Top band: ANSWERS the research question, states CONCRETE limitations (and ethics where relevant),
%   proposes SPECIFIC future work.
% Limitations to cover: 8 days only (one training weekend day, so no weekend effect: FB3); no weather,
%   traffic or courier data at order time; coordinates shifted for privacy; one city platform.
% Ethics (Lec 1): tighter promises pressure couriers; padded promises shift risk to customers;
%   do some districts get systematically worse ETAs?
% Future work (Lec 2, monitoring): drift over longer periods, retraining, input monitoring.
\section{Conclusion}

\textbf{Answer to the research question.} Using only information known when an order is placed,
gradient boosting predicts delivery time with an MAE of 6.49 minutes on unseen days: 31\% better than
a median predictor, which meets our 20\% target. With its padding removed, the platform's promise is as
accurate as our model or slightly better (6.36--6.48), so base order-time data reach the platform's own
level, not beyond it. The three model families end within about 0.3 minutes of each other, so the limit
is the information available at ordering, not the choice of model. Distance drives delivery time far
more than anything else; recent district load, which our proposal expected to be a top driver, was the
weakest. Information the platform has in real time lowers the cross-validated MAE from 6.25 to 5.92
minutes, yet closes only a fifth of the gap to a model that sees after-order times. The business
criterion was not met: at the same average promised time, our 80th-percentile promises are late slightly
more often than the platform's, which can still react after an order is placed and whose promise shaped
the outcomes we score.

\textbf{Limitations.} The data cover eight days, with a single weekend day for training, so weekend
effects could not be modelled; the Sunday test day was indeed harder than the Monday (MAE 6.77 vs 6.17),
and eight days cannot show longer-term drift. Weather and traffic are not in the dataset, rider
availability can only be approximated from past deliveries, and coordinates are shifted for privacy, so
distances are approximate. The split follows the dataset's day field, the local arrival date, so 16 test
orders were placed before the last training order. All results come from one platform in one city, over
one week.

\textbf{Ethics.} Errors are not spread evenly: test MAE ranges from 4.4 to 9.0 minutes across districts,
so customers in some areas would get systematically less reliable estimates. Promise length also moves
risk between groups: shorter promises put pressure on couriers to ride faster, while padded promises
shift the cost of uncertainty to customers. A deployed system should report errors by district and be
judged on courier workload as well as late rate.

\textbf{Future work.} (1) Confirm the extended features' cross-validated gain on new, unseen days, then
add weather and order size. (2) Combine the estimate with post-order updates, which is
where the platform's advantage appears to come from. (3) Before deployment, test on several weeks to
measure drift, retrain on a schedule, and monitor input distributions for change.
